\documentclass{article}
\usepackage{../header}
\title{Fluid Dynamics II}
\author{Notes made by Finley Cooper}

\begin{document}
  \maketitle
  \newpage
  \tableofcontents
  \newpage
\section{Introduction}
The Part IB Fluid Dynamics course mainly covers inviscid flow and irrotational flow including potential flow. The course covered an introduction to viscious flow for the very simple case where $ \mathbf u = (u(y,t), 0) $ in 2D parallel flow.
\par
This course is about the effects of viscosity and answers the following very important questions
\begin{enumerate}
	\item How does viscosity affect the equations for a general flow?
	\item What happens if viscosity dominates over the inertia of the fluid?
	\item What happens if the flow is nearly paralllel?
	\item What happens if viscosity is small compared to inertia?
	\item What about stability?
\end{enumerate}
\section{Equations for a Newtonian Viscious Flow}
\subsection{Dynamic Viscosity}
Recall from Part IB Fluid Dynamics we have Couette flow where a plate at rest and a plate with velocity $ U $ are parallel with a space $ h $ between them filled with a viscous fluid. For a simple fluid we find that the force required to move the top plate $ F $ is proportional to $ \frac {uA}h $. Hence the shear stress is
\[
  \tau_s = \frac FA = \mu\frac Uh
\]
where $ \mu $ is the dynamic viscosity. The shear stress is due to the relative motion between each layer of fluid and its neighbours above and below.
\subsection{Rate of strain tensor}
Consider the velocity near a fixed point $ \mathbf 0 $ \textit{wlog}. Taylor expand $ \mathbf u $ so we have the multivaraible Taylor expansion
\[
	u_i(\mathbf x) = u_i(0) + x_j\left.\frac{\partial u_i}{\partial x_j} \right|_{\mathbf 0} + \frac 12 x_jx_k\left.\frac{\partial^2 u_i}{\partial x_j\partial x_k}\right|_{\mathbf 0} + \cdots
\]
We drop the quadratic terms and decompose $ \frac{\partial u_i}{\partial x_j} $ into symmetric and antisymmetric parts so
\begin{align*}
	\frac{\partial u_i}{\partial x_j} &= \frac 12 \left(\frac{\partial u_i}{\partial x_j} + \frac{\partial u_j}{\partial x_i}\right) + \frac 12 \left(\frac{\partial u_i}{\partial x_j}-\frac{\partial u_j}{\partial x_i}\right) \\
					  &= e_{ij} + \frac 12 \varepsilon_{jik} \boldsymbol \omega_k
\end{align*}
where $ e_{ij} $ is the \textit{rate of strain tensor} and $ \boldsymbol\omega = \nabla\times \mathbf u  $ is the vorticity. Hence we have that
\[
  \mathbf u (\mathbf x) \approx \mathbf u (\mathbf 0) + \mathbf e \cdot \mathbf x + \frac 12 \boldsymbol \omega \times \mathbf x.
\]
We have that $ e_{ij}= e_{ji} $ by definition so $ \mathbf e $ is real symmetic and hence has real orthogonal eigenvectors which are the principle axis of strain. We have that $ e_{ii} = \nabla\cdot\mathbf u = 0 $ so the eigenvalues sum to $ 0 $. The eigenvalues are called the principle rates of strain.
\end{document}
